% Source: book pp. 354--369 (PDF 368--383); no solutions.
% \exnote: Axler's small italic comment under an exercise
\section{Determinants}

\subsection{Defining the Determinant}

The next definition will lead us to a clean, beautiful, basis-free definition
of the determinant of an operator.

\begin{definition}[$\alpha_T$]\label{ax:9.40}
Suppose that $m$ is a positive integer and $T\in\Lin(V)$. For
$\alpha\in V_{\mathrm{alt}}^{(m)}$, define $\alpha_T\in V_{\mathrm{alt}}^{(m)}$
by
\[ \alpha_T(v_1,\dots,v_m) = \alpha(Tv_1,\dots,Tv_m) \]
for each list $v_1,\dots,v_m$ of vectors in $V$.
\end{definition}

Suppose $T\in\Lin(V)$. If $\alpha\in V_{\mathrm{alt}}^{(m)}$ and
$v_1,\dots,v_m$ is a list of vectors in $V$ with $v_j=v_k$ for some $j\ne k$,
then $Tv_j=Tv_k$, which implies that
$\alpha_T(v_1,\dots,v_m)=\alpha(Tv_1,\dots,Tv_m)=0$. Thus the function
$\alpha\mapsto\alpha_T$ is a linear map of $V_{\mathrm{alt}}^{(m)}$ to itself.

\begingroup\emergencystretch=2em
We know that $\dim V_{\mathrm{alt}}^{(\dim V)}=1$ (see \ref{ax:9.37}). Every
linear map from a one-dimensional vector space to itself is multiplication by
some unique scalar. For the linear map $\alpha\mapsto\alpha_T$, we now define
$\det T$ to be that scalar.\par\endgroup

\begin{definition}[Determinant of an operator, $\det T$]\label{ax:9.41}
Suppose $T\in\Lin(V)$. The \emph{determinant} of $T$, denoted by $\det T$, is
defined to be the unique number in $\F$ such that
\[ \alpha_T = (\det T)\,\alpha \]
for all $\alpha\in V_{\mathrm{alt}}^{(\dim V)}$.
\end{definition}

\begin{example}[Determinants of operators]\label{ax:9.42}
Let $n=\dim V$.
\begin{itemize}
\item If $I$ is the identity operator on $V$, then $\alpha_I=\alpha$ for all
  $\alpha\in V_{\mathrm{alt}}^{(n)}$. Thus $\det I=1$.
\item More generally, if $\lambda\in\F$, then $\alpha_{\lambda I}=\lambda^n\alpha$
  for all $\alpha\in V_{\mathrm{alt}}^{(n)}$. Thus $\det(\lambda I)=\lambda^n$.
\item Still more generally, if $T\in\Lin(V)$ and $\lambda\in\F$, then
  $\alpha_{\lambda T}=\lambda^n\alpha_T=\lambda^n(\det T)\alpha$ for all
  $\alpha\in V_{\mathrm{alt}}^{(n)}$. Thus $\det(\lambda T)=\lambda^n\det T$.
\item Suppose $T\in\Lin(V)$ and there is a basis $e_1,\dots,e_n$ of $V$
  consisting of eigenvectors of $T$, with corresponding eigenvalues
  $\lambda_1,\dots,\lambda_n$. If $\alpha\in V_{\mathrm{alt}}^{(n)}$, then
  \[ \alpha_T(e_1,\dots,e_n) = \alpha(\lambda_1e_1,\dots,\lambda_ne_n)
     = (\lambda_1\cdots\lambda_n)\alpha(e_1,\dots,e_n). \]
  If $\alpha\ne0$, then \ref{ax:9.39} implies $\alpha(e_1,\dots,e_n)\ne0$.
  Thus the equation above implies
  \[ \det T = \lambda_1\cdots\lambda_n. \]
\end{itemize}
\end{example}

Our next task is to define and give a formula for the determinant of a square
matrix. To do this, we associate with each square matrix an operator and then
define the determinant of the matrix to be the determinant of the associated
operator.

\begin{definition}[Determinant of a matrix, $\det A$]\label{ax:9.43}
Suppose that $n$ is a positive integer and $A$ is an $n$-by-$n$ square matrix
with entries in $\F$. Let $T\in\Lin(\F^n)$ be the operator whose matrix with
respect to the standard basis of $\F^n$ equals $A$. The \emph{determinant} of
$A$, denoted by $\det A$, is defined by $\det A=\det T$.
\end{definition}

\begin{example}[Determinants of matrices]\label{ax:9.44}\leavevmode
\begin{itemize}
\item If $I$ is the $n$-by-$n$ identity matrix, then the corresponding operator
  on $\F^n$ is the identity operator $I$ on $\F^n$. Thus the first bullet point
  of \ref{ax:9.42} implies that the determinant of the identity matrix is $1$.
\item Suppose $A$ is a diagonal matrix with $\lambda_1,\dots,\lambda_n$ on the
  diagonal. Then the corresponding operator on $\F^n$ has the standard basis
  of $\F^n$ as eigenvectors, with eigenvalues $\lambda_1,\dots,\lambda_n$. Thus
  the last bullet point of \ref{ax:9.42} implies that
  $\det A=\lambda_1\cdots\lambda_n$.
\end{itemize}
\end{example}

For the next result, think of each list $v_1,\dots,v_n$ of $n$ vectors in
$\F^n$ as a list of $n$-by-$1$ column vectors. The notation
$\begin{pmatrix} v_1 & \cdots & v_n \end{pmatrix}$ then denotes the
$n$-by-$n$ square matrix whose $k^{\text{th}}$ column is $v_k$ for each
$k=1,\dots,n$.

\begin{result}[Determinant is an alternating multilinear form]\label{ax:9.45}
Suppose that $n$ is a positive integer. The map that takes a list
$v_1,\dots,v_n$ of vectors in $\F^n$ to
$\det\begin{pmatrix} v_1 & \cdots & v_n \end{pmatrix}$ is an alternating
$n$-linear form on $\F^n$.
\end{result}

\begin{proof}
Let $e_1,\dots,e_n$ be the standard basis of $\F^n$ and suppose
$v_1,\dots,v_n$ is a list of vectors in $\F^n$. Let $T\in\Lin(\F^n)$ be the
operator such that $Te_k=v_k$ for $k=1,\dots,n$. Thus $T$ is the operator whose
matrix with respect to $e_1,\dots,e_n$ is
$\begin{pmatrix} v_1 & \cdots & v_n \end{pmatrix}$. Hence
$\det\begin{pmatrix} v_1 & \cdots & v_n \end{pmatrix}=\det T$, by definition
of the determinant of a matrix.

Let $\alpha$ be an alternating $n$-linear form on $\F^n$ such that
$\alpha(e_1,\dots,e_n)=1$. Then
\begin{align*}
\det\begin{pmatrix} v_1 & \cdots & v_n \end{pmatrix} &= \det T\\
  &= (\det T)\,\alpha(e_1,\dots,e_n)\\
  &= \alpha(Te_1,\dots,Te_n)\\
  &= \alpha(v_1,\dots,v_n),
\end{align*}
where the third line follows from the definition of the determinant of an
operator. The equation above shows that the map that takes a list of vectors
$v_1,\dots,v_n$ in $\F^n$ to
$\det\begin{pmatrix} v_1 & \cdots & v_n \end{pmatrix}$ is the alternating
$n$-linear form $\alpha$ on $\F^n$.
\end{proof}

The previous result has several important consequences. For example, it
immediately implies that a matrix with two identical columns has
determinant~$0$. We will come back to other consequences later, but for now we
want to give a formula for the determinant of a square matrix. Recall that if
$A$ is a matrix, then $A_{j,k}$ denotes the entry in row $j$, column $k$ of
$A$.

\begin{result}[Formula for determinant of a matrix]\label{ax:9.46}
Suppose that $n$ is a positive integer and $A$ is an $n$-by-$n$ square matrix.
Then
\[ \det A = \sum_{(j_1,\dots,j_n)\in\operatorname{perm} n}
   \bigl(\sgn(j_1,\dots,j_n)\bigr)A_{j_1,1}\cdots A_{j_n,n}. \]
\end{result}

\begin{proof}
Apply \ref{ax:9.36} with $V=\F^n$ and $e_1,\dots,e_n$ the standard basis of
$\F^n$ and $\alpha$ the alternating $n$-linear form on $\F^n$ that takes
$v_1,\dots,v_n$ to $\det\begin{pmatrix} v_1 & \cdots & v_n \end{pmatrix}$ [see
\ref{ax:9.45}]. If each $v_k$ is the $k^{\text{th}}$ column of $A$, then each
$b_{j,k}$ in \ref{ax:9.36} equals $A_{j,k}$. Finally,
\[ \alpha(e_1,\dots,e_n) = \det\begin{pmatrix} e_1 & \cdots & e_n \end{pmatrix}
   = \det I = 1. \]
Thus the formula in \ref{ax:9.36} becomes the formula stated in this result.
\end{proof}

\begin{example}[Explicit formula for determinant]\label{ax:9.47}\leavevmode
\begin{itemize}
\item If $A$ is a $2$-by-$2$ matrix, then the formula in \ref{ax:9.46} becomes
  \[ \det A = A_{1,1}A_{2,2} - A_{2,1}A_{1,2}. \]
\item If $A$ is a $3$-by-$3$ matrix, then the formula in \ref{ax:9.46} becomes
  \begin{align*}
  \det A = {}&A_{1,1}A_{2,2}A_{3,3} - A_{2,1}A_{1,2}A_{3,3}
             - A_{3,1}A_{2,2}A_{1,3}\\
             &- A_{1,1}A_{3,2}A_{2,3} + A_{3,1}A_{1,2}A_{2,3}
             + A_{2,1}A_{3,2}A_{1,3}.
  \end{align*}
\end{itemize}
\end{example}

The sum in the formula in \ref{ax:9.46} contains $n!$ terms. Because $n!$ grows
rapidly as $n$ increases, the formula in \ref{ax:9.46} is not a viable method
to evaluate determinants even for moderately sized $n$. For example, $10!$ is
over three million, and $100!$ is approximately $10^{158}$, leading to a sum
that the fastest computer cannot evaluate. We will soon see some results that
lead to faster evaluations of determinants than direct use of the sum in
\ref{ax:9.46}.

\begin{result}[Determinant of upper-triangular matrix]\label{ax:9.48}
Suppose that $A$ is an upper-triangular matrix with
$\lambda_1,\dots,\lambda_n$ on the diagonal. Then
$\det A=\lambda_1\cdots\lambda_n$.
\end{result}

\begin{proof}
If $(j_1,\dots,j_n)\in\operatorname{perm} n$ with
$(j_1,\dots,j_n)\ne(1,\dots,n)$, then $j_k>k$ for some $k\in\{1,\dots,n\}$,
which implies that $A_{j_k,k}=0$. Thus the only permutation that can make a
nonzero contribution to the sum in \ref{ax:9.46} is the permutation
$(1,\dots,n)$. Because $A_{k,k}=\lambda_k$ for each $k=1,\dots,n$, this implies
that $\det A=\lambda_1\cdots\lambda_n$.
\end{proof}

\subsection{Properties of Determinants}

Our definition of the determinant leads to the following magical proof that
the determinant is multiplicative.

\begin{result}[Determinant is multiplicative]\label{ax:9.49}\leavevmode
\begin{parts}
\item Suppose $S,T\in\Lin(V)$. Then $\det(ST)=(\det S)(\det T)$.
\item Suppose $A$ and $B$ are square matrices of the same size. Then
  \[ \det(AB) = (\det A)(\det B) \]
\end{parts}
\end{result}

\begin{proof}\leavevmode
\begin{parts}
\item Let $n=\dim V$. Suppose $\alpha\in V_{\mathrm{alt}}^{(n)}$ and
  $v_1,\dots,v_n\in V$. Then
  \begin{align*}
  \alpha_{ST}(v_1,\dots,v_n) &= \alpha(STv_1,\dots,STv_n)\\
    &= (\det S)\alpha(Tv_1,\dots,Tv_n)\\
    &= (\det S)(\det T)\alpha(v_1,\dots,v_n),
  \end{align*}
  where the first equation follows from the definition of $\alpha_{ST}$, the
  second equation follows from the definition of $\det S$, and the third
  equation follows from the definition of $\det T$. The equation above implies
  that $\det(ST)=(\det S)(\det T)$.
\item Let $S,T\in\Lin(\F^n)$ be such that $\Mat(S)=A$ and $\Mat(T)=B$, where
  all matrices of operators in this proof are with respect to the standard
  basis of $\F^n$. Then $\Mat(ST)=\Mat(S)\Mat(T)=AB$ (see \ref{ax:3.43}). Thus
  \[ \det(AB) = \det(ST) = (\det S)(\det T) = (\det A)(\det B), \]
  where the second equality comes from the result in (a). \qedhere
\end{parts}
\end{proof}

The determinant of an operator determines whether the operator is invertible.

\begin{result}[Invertible $\iff$ nonzero determinant]\label{ax:9.50}
An operator $T\in\Lin(V)$ is invertible if and only if $\det T\ne0$.
Furthermore, if $T$ is invertible, then $\det(T^{-1})=1/(\det T)$.
\end{result}

\begin{proof}
First suppose $T$ is invertible. Thus $TT^{-1}=I$. Now \ref{ax:9.49} implies
that
\[ 1 = \det I = \det(TT^{-1}) = (\det T)\bigl(\det(T^{-1})\bigr). \]
Hence $\det T\ne0$ and $\det(T^{-1})$ is the multiplicative inverse of
$\det T$.

To prove the other direction, now suppose $\det T\ne0$. Suppose $v\in V$ and
$v\ne0$. Let $v,e_2,\dots,e_n$ be a basis of $V$ and let
$\alpha\in V_{\mathrm{alt}}^{(n)}$ be such that $\alpha\ne0$. Then
$\alpha(v,e_2,\dots,e_n)\ne0$ (by \ref{ax:9.39}). Now
\[ \alpha(Tv,Te_2,\dots,Te_n) = (\det T)\alpha(v,e_2,\dots,e_n) \ne 0. \]
Thus $Tv\ne0$. Hence $T$ is invertible.
\end{proof}

An $n$-by-$n$ matrix $A$ is invertible (see \ref{ax:3.80} for the definition of
an invertible matrix) if and only if the operator on $\F^n$ associated with
$A$ (via the standard basis of $\F^n$) is invertible. Thus the previous result
shows that a square matrix $A$ is invertible if and only if $\det A\ne0$.

\begin{result}[Eigenvalues and determinants]\label{ax:9.51}
Suppose $T\in\Lin(V)$ and $\lambda\in\F$. Then $\lambda$ is an eigenvalue of
$T$ if and only if $\det(\lambda I-T)=0$.
\end{result}

\begin{proof}
The number $\lambda$ is an eigenvalue of $T$ if and only if $T-\lambda I$ is
not invertible (see \ref{ax:5.7}), which happens if and only if $\lambda I-T$
is not invertible, which happens if and only if $\det(\lambda I-T)=0$ (by
\ref{ax:9.50}).
\end{proof}

Suppose $T\in\Lin(V)$ and $S\colon W\to V$ is an invertible linear map. To
prove that $\det(S^{-1}TS)=\det T$, we could try to use \ref{ax:9.49} and
\ref{ax:9.50}, writing
\begin{align*}
\det(S^{-1}TS) &= (\det S^{-1})(\det T)(\det S)\\
               &= \det T.
\end{align*}
That proof works if $W=V$, but if $W\ne V$ then it makes no sense because the
determinant is defined only for linear maps from a vector space to itself, and
$S$ maps $W$ to $V$, making $\det S$ undefined. The proof given below works
around this issue and is valid when $W\ne V$.

\begin{result}[Determinant is a similarity invariant]\label{ax:9.52}
Suppose $T\in\Lin(V)$ and $S\colon W\to V$ is an invertible linear map. Then
\[ \det(S^{-1}TS) = \det T. \]
\end{result}

\begin{proof}
Let $n=\dim W=\dim V$. Suppose $\tau\in W_{\mathrm{alt}}^{(n)}$. Define
$\alpha\in V_{\mathrm{alt}}^{(n)}$ by
\[ \alpha(v_1,\dots,v_n) = \tau(S^{-1}v_1,\dots,S^{-1}v_n) \]
for $v_1,\dots,v_n\in V$. Suppose $w_1,\dots,w_n\in W$. Then
\begin{align*}
\tau_{S^{-1}TS}(w_1,\dots,w_n) &= \tau(S^{-1}TSw_1,\dots,S^{-1}TSw_n)\\
  &= \alpha(TSw_1,\dots,TSw_n)\\
  &= \alpha_T(Sw_1,\dots,Sw_n)\\
  &= (\det T)\alpha(Sw_1,\dots,Sw_n)\\
  &= (\det T)\tau(w_1,\dots,w_n).
\end{align*}
The equation above and the definition of the determinant of the operator
$S^{-1}TS$ imply that $\det(S^{-1}TS)=\det T$.
\end{proof}

For the special case in which $V=\F^n$ and $e_1,\dots,e_n$ is the standard
basis of $\F^n$, the next result is true by the definition of the determinant
of a matrix. The left side of the equation in the next result does not depend
on a choice of basis, which means that the right side is independent of the
choice of basis.

\begin{result}[Determinant of operator equals determinant of its matrix]\label{ax:9.53}
Suppose $T\in\Lin(V)$ and $e_1,\dots,e_n$ is a basis of $V$. Then
\[ \det T = \det\Mat\bigl(T,(e_1,\dots,e_n)\bigr). \]
\end{result}

\begin{proof}
Let $f_1,\dots,f_n$ be the standard basis of $\F^n$. Let $S\colon\F^n\to V$ be
the linear map such that $Sf_k=e_k$ for each $k=1,\dots,n$. Thus
$\Mat\bigl(S,(f_1,\dots,f_n),(e_1,\dots,e_n)\bigr)$ and
$\Mat\bigl(S^{-1},(e_1,\dots,e_n),(f_1,\dots,f_n)\bigr)$ both equal the
$n$-by-$n$ identity matrix. Hence
\refstepcounter{theorem}%
\[ \Mat\bigl(S^{-1}TS,(f_1,\dots,f_n)\bigr) = \Mat\bigl(T,(e_1,\dots,e_n)\bigr),
   \tag*{\thetheorem}\label{ax:9.54} \]
as follows from two applications of \ref{ax:3.43}. Thus
\begin{align*}
\det T &= \det(S^{-1}TS)\\
       &= \det\Mat\bigl(S^{-1}TS,(f_1,\dots,f_n)\bigr)\\
       &= \det\Mat\bigl(T,(e_1,\dots,e_n)\bigr),
\end{align*}
where the first line comes from \ref{ax:9.52}, the second line comes from the
definition of the determinant of a matrix, and the third line follows from
\ref{ax:9.54}.
\end{proof}

The next result gives a more intuitive way to think about determinants than the
definition or the formula in \ref{ax:9.46}. We could make the characterization
in the result below the definition of the determinant of an operator on a
finite-dimensional complex vector space, with the current definition then
becoming a consequence of that definition.

\begin{result}[If $\F=\C$, then determinant equals product of eigenvalues]\label{ax:9.55}
Suppose $\F=\C$ and $T\in\Lin(V)$. Then $\det T$ equals the product of the
eigenvalues of $T$, with each eigenvalue included as many times as its
multiplicity.
\end{result}

\begin{proof}
There is a basis of $V$ with respect to which $T$ has an upper-triangular
matrix with the diagonal entries of the matrix consisting of the eigenvalues of
$T$, with each eigenvalue included as many times as its multiplicity---see
\ref{ax:8.37}. Thus \ref{ax:9.53} and \ref{ax:9.48} imply that $\det T$ equals
the product of the eigenvalues of $T$, with each eigenvalue included as many
times as its multiplicity.
\end{proof}

As the next result shows, the determinant interacts nicely with the transpose
of a square matrix, with the dual of an operator, and with the adjoint of an
operator on an inner product space.

\begin{result}[Determinant of transpose, dual, or adjoint]\label{ax:9.56}\leavevmode
\begin{parts}
\item Suppose $A$ is a square matrix. Then $\det A^{\mathrm{t}}=\det A$.
\item Suppose $T\in\Lin(V)$. Then $\det T'=\det T$.
\item Suppose $V$ is an inner product space and $T\in\Lin(V)$. Then
  \[ \det(T^*) = \overline{\det T}. \]
\end{parts}
\end{result}

\begin{proof}\leavevmode
\begin{parts}
\item Let $n$ be a positive integer. Define $\alpha\colon(\F^n)^n\to\F$ by
  \[ \alpha(v_1,\dots,v_n)
     = \det\Bigl(\begin{pmatrix} v_1 & \cdots & v_n \end{pmatrix}^{\mathrm{t}}\Bigr) \]
  for all $v_1,\dots,v_n\in\F^n$. The formula in \ref{ax:9.46} for the
  determinant of a matrix shows that $\alpha$ is an $n$-linear form on $\F^n$.

  Suppose $v_1,\dots,v_n\in\F^n$ and $v_j=v_k$ for some $j\ne k$. If $B$ is an
  $n$-by-$n$ matrix, then
  $\begin{pmatrix} v_1 & \cdots & v_n \end{pmatrix}^{\mathrm{t}}B$ cannot equal
  the identity matrix because row $j$ and row $k$ of
  $\begin{pmatrix} v_1 & \cdots & v_n \end{pmatrix}^{\mathrm{t}}B$ are equal.
  Thus $\begin{pmatrix} v_1 & \cdots & v_n \end{pmatrix}^{\mathrm{t}}$ is not
  invertible, which implies that $\alpha(v_1,\dots,v_n)=0$. Hence $\alpha$ is
  an alternating $n$-linear form on $\F^n$.

  Note that $\alpha$ applied to the standard basis of $\F^n$ equals $1$.
  Because the vector space of alternating $n$-linear forms on $\F^n$ has
  dimension one (by \ref{ax:9.37}), this implies that $\alpha$ is the
  determinant function. Thus (a) holds.
\item The equation $\det T'=\det T$ follows from (a) and \ref{ax:9.53} and
  \ref{ax:3.132}.
\item Pick an orthonormal basis of $V$. The matrix of $T^*$ with respect to
  that basis is the conjugate transpose of the matrix of $T$ with respect to
  that basis (by \ref{ax:7.9}). Thus \ref{ax:9.53}, \ref{ax:9.46}, and (a)
  imply that $\det(T^*)=\overline{\det T}$. \qedhere
\end{parts}
\end{proof}

\begin{result}[Helpful results in evaluating determinants]\label{ax:9.57}\leavevmode
\begin{parts}
\item If either two columns or two rows of a square matrix are equal, then the
  determinant of the matrix equals $0$.
\item Suppose $A$ is a square matrix and $B$ is the matrix obtained from $A$ by
  swapping either two columns or two rows. Then $\det A=-\det B$.
\item If one column or one row of a square matrix is multiplied by a scalar,
  then the value of the determinant is multiplied by the same scalar.
\item If a scalar multiple of one column of a square matrix is added to another
  column, then the value of the determinant is unchanged.
\item If a scalar multiple of one row of a square matrix is added to another
  row, then the value of the determinant is unchanged.
\end{parts}
\end{result}

\begin{proof}
All the assertions in this result follow from the result that the maps
$v_1,\dots,v_n\mapsto\det\begin{pmatrix} v_1 & \cdots & v_n \end{pmatrix}$ and
$v_1,\dots,v_n\mapsto
\det\begin{pmatrix} v_1 & \cdots & v_n \end{pmatrix}^{\mathrm{t}}$ are both
alternating $n$-linear forms on $\F^n$ [see \ref{ax:9.45} and
\ref{ax:9.56}(a)].

For example, to prove (d) suppose $v_1,\dots,v_n\in\F^n$ and $c\in\F$. Then
\begin{align*}
&\det\begin{pmatrix} v_1+cv_2 & v_2 & \cdots & v_n \end{pmatrix}\\
&\qquad= \det\begin{pmatrix} v_1 & v_2 & \cdots & v_n \end{pmatrix}
  + c\det\begin{pmatrix} v_2 & v_2 & v_3 & \cdots & v_n \end{pmatrix}\\
&\qquad= \det\begin{pmatrix} v_1 & v_2 & \cdots & v_n \end{pmatrix},
\end{align*}
where the first equation follows from the multilinearity property and the
second equation follows from the alternating property. The equation above
shows that adding a multiple of the second column to the first column does not
change the value of the determinant. The same conclusion holds for any two
columns. Thus (d) holds.

The proof of (e) follows from (d) and from \ref{ax:9.56}(a). The proofs of (a),
(b), and (c) use similar tools and are left to the reader.
\end{proof}

For matrices whose entries are concrete numbers, the result above leads to a
much faster way to evaluate the determinant than direct application of the
formula in \ref{ax:9.46}. Specifically, apply the Gaussian elimination
procedure of swapping rows [by \ref{ax:9.57}(b), this changes the determinant
by a factor of $-1$], multiplying a row by a nonzero constant [by
\ref{ax:9.57}(c), this changes the determinant by the same constant], and
adding a multiple of one row to another row [by \ref{ax:9.57}(e), this does not
change the determinant] to produce an upper-triangular matrix, whose
determinant is the product of the diagonal entries (by \ref{ax:9.48}). If your
software keeps track of the number of row swaps and of the constants used when
multiplying a row by a constant, then the determinant of the original matrix
can be computed.

Because a number $\lambda\in\F$ is an eigenvalue of an operator $T\in\Lin(V)$
if and only if $\det(\lambda I-T)=0$ (by \ref{ax:9.51}), you may be tempted to
think that one way to find eigenvalues quickly is to choose a basis of $V$, let
$A=\Mat(T)$, evaluate $\det(\lambda I-A)$, and then solve the equation
$\det(\lambda I-A)=0$ for $\lambda$. However, that procedure is rarely
efficient, except when $\dim V=2$ (or when $\dim V$ equals $3$ or $4$ if you
are willing to use the cubic or quartic formulas). One problem is that the
procedure described in the paragraph above for evaluating a determinant does
not work when the matrix includes a symbol (such as the $\lambda$ in
$\lambda I-A$). This problem arises because decisions need to be made in the
Gaussian elimination procedure about whether certain quantities equal $0$, and
those decisions become complicated in expressions involving a symbol
$\lambda$.

Recall that an operator on a finite-dimensional inner product space is unitary
if it preserves norms (see \ref{ax:7.51} and the paragraph following it). Every
eigenvalue of a unitary operator has absolute value $1$ (by \ref{ax:7.54}).
Thus the product of the eigenvalues of a unitary operator has absolute
value~$1$. Hence (at least in the case $\F=\C$) the determinant of a unitary
operator has absolute value $1$ (by \ref{ax:9.55}). The next result gives a
proof that works without the assumption that $\F=\C$.

\begin{result}[Every unitary operator has determinant with absolute value $1$]\label{ax:9.58}
Suppose $V$ is an inner product space and $S\in\Lin(V)$ is a unitary operator.
Then $\abs{\det S}=1$.
\end{result}

\begin{proof}
Because $S$ is unitary, $I=S^*S$ (see \ref{ax:7.53}). Thus
\[ 1 = \det(S^*S) = (\det S^*)(\det S) = (\overline{\det S})(\det S)
   = \abs{\det S}^2, \]
where the second equality comes from \ref{ax:9.49}(a) and the third equality
comes from \ref{ax:9.56}(c). The equation above implies that
$\abs{\det S}=1$.
\end{proof}

The determinant of a positive operator on an inner product space meshes well
with the analogy that such operators correspond to the nonnegative real
numbers.

\begin{result}[Every positive operator has nonnegative determinant]\label{ax:9.59}
Suppose $V$ is an inner product space and $T\in\Lin(V)$ is a positive
operator. Then $\det T\ge0$.
\end{result}

\begin{proof}
By the spectral theorem (\ref{ax:7.29} or \ref{ax:7.31}), $V$ has an
orthonormal basis consisting of eigenvectors of $T$. Thus by the last bullet
point of \ref{ax:9.42}, $\det T$ equals a product of the eigenvalues of $T$,
possibly with repetitions. Each eigenvalue of $T$ is a nonnegative number (by
\ref{ax:7.38}). Thus we conclude that $\det T\ge0$.
\end{proof}

Suppose $V$ is an inner product space and $T\in\Lin(V)$. Recall that the list
of nonnegative square roots of the eigenvalues of $T^*T$ (each included as many
times as its multiplicity) is called the list of singular values of $T$ (see
Section~7E).

\begin{result}[$\abs{\det T}={}$product of singular values of $T$]\label{ax:9.60}
Suppose $V$ is an inner product space and $T\in\Lin(V)$. Then
\[ \abs{\det T} = \sqrt{\det(T^*T)}
   = \text{product of singular values of } T. \]
\end{result}

\begin{proof}
We have
\[ \abs{\det T}^2 = (\overline{\det T})(\det T)
   = \bigl(\det(T^*)\bigr)(\det T) = \det(T^*T), \]
where the middle equality comes from \ref{ax:9.56}(c) and the last equality
comes from \ref{ax:9.49}(a). Taking square roots of both sides of the equation
above shows that $\abs{\det T}=\sqrt{\det(T^*T)}$.

Let $s_1,\dots,s_n$ denote the list of singular values of $T$. Thus
$s_1^2,\dots,s_n^2$ is the list of eigenvalues of $T^*T$ (with appropriate
repetitions), corresponding to an orthonormal basis of $V$ consisting of
eigenvectors of $T^*T$. Hence the last bullet point of \ref{ax:9.42} implies
that
\[ \det(T^*T) = s_1^2\cdots s_n^2. \]
Thus $\abs{\det T}=s_1\cdots s_n$, as desired.
\end{proof}

An operator $T$ on a real inner product space changes volume by a factor of the
product of the singular values (by \ref{ax:7.111}). Thus the next result
follows immediately from \ref{ax:7.111} and \ref{ax:9.60}. This result explains
why the absolute value of a determinant appears in the change of variables
formula in multivariable calculus.

\begin{result}[$T$ changes volume by factor of $\abs{\det T}$]\label{ax:9.61}
Suppose $T\in\Lin(\R^n)$ and $\Omega\subseteq\R^n$. Then
\[ \operatorname{volume} T(\Omega)
   = \abs{\det T}(\operatorname{volume}\Omega). \]
\end{result}

For operators on finite-dimensional complex vector spaces, we now connect the
determinant to a polynomial that we have previously seen.

\begin{result}[If $\F=\C$, then characteristic polynomial of $T$ equals $\det(zI-T)$]\label{ax:9.62}
Suppose $\F=\C$ and $T\in\Lin(V)$. Let $\lambda_1,\dots,\lambda_m$ denote the
distinct eigenvalues of $T$, and let $d_1,\dots,d_m$ denote their
multiplicities. Then
\[ \det(zI-T) = (z-\lambda_1)^{d_1}\cdots(z-\lambda_m)^{d_m}. \]
\end{result}

\begin{proof}
There exists a basis of $V$ with respect to which $T$ has an upper-triangular
matrix with each $\lambda_k$ appearing on the diagonal exactly $d_k$ times (by
\ref{ax:8.37}). With respect to this basis, $zI-T$ has an upper-triangular
matrix with $z-\lambda_k$ appearing on the diagonal exactly $d_k$ times for
each $k$. Thus \ref{ax:9.48} gives the desired equation.
\end{proof}

Suppose $\F=\C$ and $T\in\Lin(V)$. The characteristic polynomial of $T$ was
defined in \ref{ax:8.26} as the polynomial on the right side of the equation in
\ref{ax:9.62}. We did not previously define the characteristic polynomial of an
operator on a finite-dimensional real vector space because such operators may
have no eigenvalues, making a definition using the right side of the equation
in \ref{ax:9.62} inappropriate.

We now present a new definition of the characteristic polynomial, motivated by
\ref{ax:9.62}. This new definition is valid for both real and complex vector
spaces. The equation in \ref{ax:9.62} shows that this new definition is
equivalent to our previous definition when $\F=\C$ (\ref{ax:8.26}).

\begin{definition}[Characteristic polynomial]\label{ax:9.63}
Suppose $T\in\Lin(V)$. The polynomial defined by
\[ z\mapsto\det(zI-T) \]
is called the \emph{characteristic polynomial} of $T$.
\end{definition}

The formula in \ref{ax:9.46} shows that the characteristic polynomial of an
operator $T\in\Lin(V)$ is a monic polynomial of degree $\dim V$. The zeros in
$\F$ of the characteristic polynomial of $T$ are exactly the eigenvalues of $T$
(by \ref{ax:9.51}).

Previously we proved the Cayley--Hamilton theorem (\ref{ax:8.29}) in the
complex case. Now we can extend that result to operators on real vector spaces.

\begin{result}[Cayley--Hamilton theorem]\label{ax:9.64}
Suppose $T\in\Lin(V)$ and $q$ is the characteristic polynomial of $T$. Then
$q(T)=0$.
\end{result}

\begin{proof}
If $\F=\C$, then the equation $q(T)=0$ follows from \ref{ax:9.62} and
\ref{ax:8.29}.

Now suppose $\F=\R$. Fix a basis of $V$, and let $A$ be the matrix of $T$ with
respect to this basis. Let $S$ be the operator on $\C^{\dim V}$ such that the
matrix of $S$ (with respect to the standard basis of $\C^{\dim V}$) is $A$. For
all $z\in\R$ we have
\[ q(z) = \det(zI-T) = \det(zI-A) = \det(zI-S). \]
Thus $q$ is the characteristic polynomial of $S$. The case $\F=\C$ (first
sentence of this proof) now implies that $0=q(S)=q(A)=q(T)$.
\end{proof}

The Cayley--Hamilton theorem (\ref{ax:9.64}) implies that the characteristic
polynomial of an operator $T\in\Lin(V)$ is a polynomial multiple of the minimal
polynomial of $T$ (by \ref{ax:5.29}). Thus if the degree of the minimal
polynomial of $T$ equals $\dim V$, then the characteristic polynomial of $T$
equals the minimal polynomial of $T$. This happens for a very large percentage
of operators, including over $99.999\%$ of $4$-by-$4$ matrices with integer
entries in $[-100,100]$ (see the paragraph following \ref{ax:5.25}).

The last sentence in our next result was previously proved in the complex case
(see \ref{ax:8.54}). Now we can give a proof that works on both real and
complex vector spaces.

\begin{result}[Characteristic polynomial, trace, and determinant]\label{ax:9.65}
Suppose $T\in\Lin(V)$. Let $n=\dim V$. Then the characteristic polynomial of
$T$ can be written as
\[ z^n - (\tr T)z^{n-1} + \cdots + (-1)^n(\det T). \]
\end{result}

\begin{proof}
The constant term of a polynomial function of $z$ is the value of the
polynomial when $z=0$. Thus the constant term of the characteristic polynomial
of $T$ equals $\det(-T)$, which equals $(-1)^n\det T$ (by the third bullet
point of \ref{ax:9.42}).

Fix a basis of $V$, and let $A$ be the matrix of $T$ with respect to this
basis. The matrix of $zI-T$ with respect to this basis is $zI-A$. The term
coming from the identity permutation $\{1,\dots,n\}$ in the formula
\ref{ax:9.46} for $\det(zI-A)$ is
\[ (z-A_{1,1})\cdots(z-A_{n,n}). \]
The coefficient of $z^{n-1}$ in the expression above is
$-(A_{1,1}+\cdots+A_{n,n})$, which equals $-\tr T$. The terms in the formula
for $\det(zI-A)$ coming from other elements of $\operatorname{perm} n$ contain
at most $n-2$ factors of the form $z-A_{k,k}$ and thus do not contribute to the
coefficient of $z^{n-1}$ in the characteristic polynomial of $T$.
\end{proof}

\marginnote{The next result was proved by Jacques Hadamard (1865--1963) in
1893.}%
In the result below, think of the columns of the $n$-by-$n$ matrix $A$ as
elements of $\F^n$. The norms appearing below then arise from the standard
inner product on $\F^n$. Recall that the notation $R_{\cdot,k}$ in the proof
below means the $k^{\text{th}}$ column of the matrix $R$ (as was defined in
\ref{ax:3.44}).

\begin{result}[Hadamard's inequality]\label{ax:9.66}
Suppose $A$ is an $n$-by-$n$ matrix. Let $v_1,\dots,v_n$ denote the columns of
$A$. Then
\[ \abs{\det A} \le \prod_{k=1}^n\norm{v_k}. \]
\end{result}

\begin{proof}
If $A$ is not invertible, then $\det A=0$ and hence the desired inequality
holds in this case.

Thus assume that $A$ is invertible. The QR factorization (\ref{ax:7.58}) tells
us that there exist a unitary matrix $Q$ and an upper-triangular matrix $R$
whose diagonal contains only positive numbers such that $A=QR$. We have
\begin{align*}
\abs{\det A} &= \abs{\det Q}\,\abs{\det R}\\
  &= \abs{\det R}\\
  &= \prod_{k=1}^n R_{k,k}\\
  &\le \prod_{k=1}^n\norm{R_{\cdot,k}}\\
  &= \prod_{k=1}^n\norm{QR_{\cdot,k}}\\
  &= \prod_{k=1}^n\norm{v_k},
\end{align*}
where the first line comes from \ref{ax:9.49}(b), the second line comes from
\ref{ax:9.58}, the third line comes from \ref{ax:9.48}, and the fifth line
holds because $Q$ is an isometry.
\end{proof}

To give a geometric interpretation to Hadamard's inequality, suppose $\F=\R$.
Let $T\in\Lin(\R^n)$ be the operator such that $Te_k=v_k$ for each
$k=1,\dots,n$, where $e_1,\dots,e_n$ is the standard basis of $\R^n$. Then $T$
maps the box $P(e_1,\dots,e_n)$ onto the parallelepiped $P(v_1,\dots,v_n)$
[see \ref{ax:7.102} and \ref{ax:7.105} for a review of this notation and
terminology]. Because the box $P(e_1,\dots,e_n)$ has volume $1$, this implies
(by \ref{ax:9.61}) that the parallelepiped $P(v_1,\dots,v_n)$ has volume
$\abs{\det T}$, which equals $\abs{\det A}$. Thus Hadamard's inequality above
can be interpreted to say that among all parallelepipeds whose edges have
lengths $\norm{v_1},\dots,\norm{v_n}$, the ones with largest volume have
orthogonal edges $($and thus have volume
$\prod_{k=1}^n\norm{v_k})$.

For a necessary and sufficient condition for Hadamard's inequality to be an
equality, see Exercise~18.

\begingroup\emergencystretch=2em
The matrix in the next result is called the \emph{Vandermonde matrix}.
Vandermonde matrices have important applications in polynomial interpolation,
the discrete Fourier transform, and other areas of mathematics. The proof of
the next result is a nice illustration of the power of switching between
matrices and linear maps.\par\endgroup

\begin{result}[Determinant of Vandermonde matrix]\label{ax:9.67}
Suppose $n>1$ and $\beta_1,\dots,\beta_n\in\F$. Then
\[ \det\begin{pmatrix}
     1 & \beta_1 & \beta_1^2 & \cdots & \beta_1^{n-1}\\
     1 & \beta_2 & \beta_2^2 & \cdots & \beta_2^{n-1}\\
       &         & \ddots    &        &              \\
     1 & \beta_n & \beta_n^2 & \cdots & \beta_n^{n-1}
   \end{pmatrix}
   = \prod_{1\le j<k\le n}(\beta_k-\beta_j). \]
\end{result}

\begin{proof}
Let $1,z,\dots,z^{n-1}$ be the standard basis of $\Poly_{n-1}(\F)$ and let
$e_1,\dots,e_n$ denote the standard basis of $\F^n$. Define a linear map
$S\colon\Poly_{n-1}(\F)\to\F^n$ by
\[ Sp = \bigl(p(\beta_1),\dots,p(\beta_n)\bigr). \]

Let $A$ denote the Vandermonde matrix shown in the statement of this result.
Note that
\[ A = \Mat\bigl(S,(1,z,\dots,z^{n-1}),(e_1,\dots,e_n)\bigr). \]

Let $T\colon\Poly_{n-1}(\F)\to\Poly_{n-1}(\F)$ be the operator on
$\Poly_{n-1}(\F)$ such that $T1=1$ and
\[ Tz^k = (z-\beta_1)(z-\beta_2)\cdots(z-\beta_k) \]
for $k=1,\dots,n-1$. Let
$B=\Mat\bigl(T,(1,z,\dots,z^{n-1}),(1,z,\dots,z^{n-1})\bigr)$. Then $B$ is an
upper-triangular matrix all of whose diagonal entries equal $1$. Thus
$\det B=1$ (by \ref{ax:9.48}).

Let $C=\Mat\bigl(ST,(1,z,\dots,z^{n-1}),(e_1,\dots,e_n)\bigr)$. Thus $C=AB$
(by \ref{ax:3.81}), which implies that
\[ \det A = (\det A)(\det B) = \det C. \]
The definitions of $C$, $S$, and $T$ show that $C$ equals
\[ \begin{pmatrix}
   1 & 0 & 0 & \cdots & 0\\
   1 & \beta_2-\beta_1 & 0 & \cdots & 0\\
   1 & \beta_3-\beta_1 & (\beta_3-\beta_1)(\beta_3-\beta_2) & \cdots & 0\\
     &  &  & \ddots & \\
   1 & \beta_n-\beta_1 & (\beta_n-\beta_1)(\beta_n-\beta_2) & \cdots
     & (\beta_n-\beta_1)(\beta_n-\beta_2)\cdots(\beta_n-\beta_{n-1})
   \end{pmatrix}. \]
Now $\det A=\det C=\displaystyle\prod_{1\le j<k\le n}(\beta_k-\beta_j)$, where
we have used \ref{ax:9.56}(a) and \ref{ax:9.48}.
\end{proof}

% ------------------------------------------------------------------ exercises
\begin{exc}

\begin{exercise}
Prove or give a counterexample: $S,T\in\Lin(V)\implies
\det(S+T)=\det S+\det T$.
\end{exercise}

\begin{exercise}
Suppose the first column of a square matrix $A$ consists of all zeros except
possibly the first entry $A_{1,1}$. Let $B$ be the matrix obtained from $A$ by
deleting the first row and the first column of $A$. Show that
$\det A=A_{1,1}\det B$.
\end{exercise}

\begin{exercise}
Suppose $T\in\Lin(V)$ is nilpotent. Prove that $\det(I+T)=1$.
\end{exercise}

\begin{exercise}
Suppose $S\in\Lin(V)$. Prove that $S$ is unitary if and only if
$\abs{\det S}=\norm{S}=1$.
\end{exercise}

\begin{exercise}
Suppose $A$ is a block upper-triangular matrix
\[ A = \begin{pmatrix}
       A_1 &        & *\\
           & \ddots &  \\
       0   &        & A_m
       \end{pmatrix}, \]
where each $A_k$ along the diagonal is a square matrix. Prove that
\[ \det A = (\det A_1)\cdots(\det A_m). \]
\end{exercise}

\begin{exercise}
Suppose $A=\begin{pmatrix} v_1 & \cdots & v_n \end{pmatrix}$ is an $n$-by-$n$
matrix, with $v_k$ denoting the $k^{\text{th}}$ column of $A$. Show that if
$(m_1,\dots,m_n)\in\operatorname{perm} n$, then
\[ \det\begin{pmatrix} v_{m_1} & \cdots & v_{m_n} \end{pmatrix}
   = \bigl(\sgn(m_1,\dots,m_n)\bigr)\det A. \]
\end{exercise}

\begin{exercise}
Suppose $T\in\Lin(V)$ is invertible. Let $p$ denote the characteristic
polynomial of $T$ and let $q$ denote the characteristic polynomial of $T^{-1}$.
Prove that
\[ q(z) = \frac{1}{p(0)}z^{\dim V}p\Bigl(\frac{1}{z}\Bigr) \]
for all nonzero $z\in\F$.
\end{exercise}

\begin{exercise}
Suppose $T\in\Lin(V)$ is an operator with no eigenvalues (which implies that
$\F=\R$). Prove that $\det T>0$.
\end{exercise}

\begin{exercise}
Suppose that $V$ is a real vector space of even dimension, $T\in\Lin(V)$, and
$\det T<0$. Prove that $T$ has at least two distinct eigenvalues.
\end{exercise}

\begin{exercise}
Suppose $V$ is a real vector space of odd dimension and $T\in\Lin(V)$. Without
using the minimal polynomial, prove that $T$ has an eigenvalue.
\exnote{\emergencystretch=2em This result was previously proved without using determinants or the
characteristic polynomial---see \ref{ax:5.34}.}
\end{exercise}

\begin{exercise}
Prove or give a counterexample: If $\F=\R$, $T\in\Lin(V)$, and $\det T>0$, then
$T$ has a square root.
\exnote{If $\F=\C$, $T\in\Lin(V)$, and $\det T\ne0$, then $T$ has a square root
(see \ref{ax:8.41}).}
\end{exercise}

\begin{exercise}
Suppose $S,T\in\Lin(V)$ and $S$ is invertible. Define $p\colon\F\to\F$ by
\[ p(z) = \det(zS-T). \]
Prove that $p$ is a polynomial of degree $\dim V$ and that the coefficient of
$z^{\dim V}$ in this polynomial is $\det S$.
\end{exercise}

\begin{exercise}
Suppose $\F=\C$, $T\in\Lin(V)$, and $n=\dim V>2$. Let
$\lambda_1,\dots,\lambda_n$ denote the eigenvalues of $T$, with each eigenvalue
included as many times as its multiplicity.
\begin{parts}
\item Find a formula for the coefficient of $z^{n-2}$ in the characteristic
  polynomial of $T$ in terms of $\lambda_1,\dots,\lambda_n$.
\item Find a formula for the coefficient of $z$ in the characteristic
  polynomial of $T$ in terms of $\lambda_1,\dots,\lambda_n$.
\end{parts}
\end{exercise}

\begin{exercise}
Suppose $V$ is an inner product space and $T$ is a positive operator on $V$.
Prove that
\[ \det\sqrt{T} = \sqrt{\det T}. \]
\end{exercise}

\begin{exercise}
Suppose $V$ is an inner product space and $T\in\Lin(V)$. Use the polar
decomposition to give a proof that
\[ \abs{\det T} = \sqrt{\det(T^*T)} \]
that is different from the proof given earlier (see \ref{ax:9.60}).
\end{exercise}

\begin{exercise}
Suppose $T\in\Lin(V)$. Define $g\colon\F\to\F$ by $g(x)=\det(I+xT)$. Show that
$g'(0)=\tr T$.
\exnote{Look for a clean solution to this exercise, without using the explicit
but complicated formula for the determinant of a matrix.}
\end{exercise}

\begin{exercise}
Suppose $a,b,c$ are positive numbers. Find the volume of the ellipsoid
\[ \biggl\{(x,y,z)\in\R^3 : \frac{x^2}{a^2}+\frac{y^2}{b^2}+\frac{z^2}{c^2}<1
   \biggr\} \]
by finding a set $\Omega\subseteq\R^3$ whose volume you know and an operator
$T$ on $\R^3$ such that $T(\Omega)$ equals the ellipsoid above.
\end{exercise}

\begin{exercise}
Suppose that $A$ is an invertible square matrix. Prove that Hadamard's
inequality (\ref{ax:9.66}) is an equality if and only if each column of $A$ is
orthogonal to the other columns.
\end{exercise}

\begin{exercise}
Suppose $V$ is an inner product space, $e_1,\dots,e_n$ is an orthonormal basis
of $V$, and $T\in\Lin(V)$ is a positive operator.
\begin{parts}
\item Prove that $\det T\le\prod_{k=1}^n\ip{Te_k}{e_k}$.
\item Prove that if $T$ is invertible, then the inequality in (a) is an equality
  if and only if $e_k$ is an eigenvector of $T$ for each $k=1,\dots,n$.
\end{parts}
\end{exercise}

\begin{exercise}
Suppose $A$ is an $n$-by-$n$ matrix, and suppose $c$ is such that
$\abs{A_{j,k}}\le c$ for all $j,k\in\{1,\dots,n\}$. Prove that
\[ \abs{\det A} \le c^nn^{n/2}. \]
\exnote{The formula for the determinant of a matrix (\ref{ax:9.46}) shows that
$\abs{\det A}\le c^nn!$. However, the estimate given by this exercise is much
better. For example, if $c=1$ and $n=100$, then $c^nn!\approx10^{158}$, but the
estimate given by this exercise is the much smaller number $10^{100}$. If $n$
is an integer power of $2$, then the inequality above is sharp and cannot be
improved.}
\end{exercise}

\begin{exercise}
Suppose $n$ is a positive integer and $\delta\colon\C^{n,n}\to\C$ is a function
such that
\[ \delta(AB) = \delta(A)\cdot\delta(B) \]
for all $A,B\in\C^{n,n}$ and $\delta(A)$ equals the product of the diagonal
entries of $A$ for each diagonal matrix $A\in\C^{n,n}$. Prove that
\[ \delta(A) = \det A \]
for all $A\in\C^{n,n}$.
\exnote{Recall that $\C^{n,n}$ denotes the set of $n$-by-$n$ matrices with
entries in $\C$. This exercise shows that the determinant is the unique
function defined on square matrices that is multiplicative and has the desired
behavior on diagonal matrices. This result is analogous to Exercise~10 in
Section~8D, which shows that the trace is uniquely determined by its algebraic
properties.}
\end{exercise}
\end{exc}

\axquote{I find that in my own elementary lectures, I have, for pedagogical
reasons, pushed determinants more and more into the background. Too often I
have had the experience that, while the students acquired facility with the
formulas, which are so useful in abbreviating long expressions, they often
failed to gain familiarity with their \emph{meaning}, and skill in manipulation
prevented the student from going into all the details of the subject and so
gaining a mastery.}{\emph{Elementary Mathematics from an Advanced Standpoint:
Geometry}, Felix Klein}
